\scanpage{208}
\mspage{-2.11-}
% archive sheet number: 207
Therefore, since $\handwritten{\sigma}^*$ is a group action:
\[
\handwritten{\sigma}^*\handwritten{\bigl(}(\handwritten{\psi}_1\cdot b_2)^{-1},\handwritten{\psi_{\unsure{x|\chi}^*}\circ\phi}\handwritten{\bigr)}=\handwritten{\sigma}^*\handwritten{\bigl[}(b_2^{-1}\cdot\handwritten{\langle\lambda\rangle}\cdot b_2),\handwritten{\sigma}^*(b_2^{-1}\cdot\handwritten{\psi}_2,x^*)\handwritten{\bigr]}.
\tag*{(*)}
\]
But $(\handwritten{\psi}_1\cdot b_2)^{-1}$ is an $R$-valued point of $G$, and $\handwritten{\psi_{x^*}\circ\phi}$ is an $R$-valued point of $\handwritten{\mathbb{A}}^n$, therefore the term on the left is an $R$-valued point of $\handwritten{\mathbb{A}}^n$.

Now, if $f$ is a $\operatorname{Spec}(A)$-valued point of $\handwritten{\mathbb{A}}^n$, for some ring $A$, let us denote by $X_i(f)$ the element of $A$ which is its $i^{\mathrm{th}}$ coordinate. Then, recalling the definition of $\handwritten{\langle\lambda\rangle}$, and that $b_2^{-1}\cdot\handwritten{\lambda}\cdot b_2$ is diagonalized, equation (*) reads, in coordinates:
\[
X_i\handwritten{\bigl\{}\handwritten{\sigma}^*\handwritten{\bigl(}(\handwritten{\psi}_1 b_2)^{-1},\handwritten{\psi_{x^*}\circ\phi}\handwritten{\bigr)\bigr\}}=t^{\handwritten{\unsure{r|n}_i}}\cdot X_i\handwritten{\bigl\{}\handwritten{\sigma}^*(b_2^{-1}\cdot\handwritten{\psi}_2,x^*)\handwritten{\bigr\}}
\tag*{(*)$\handwritten{'}$}
\]
and, in particular:
\[
X_i\handwritten{\bigl\{}\handwritten{\sigma}^*(b_2^{-1}\cdot\handwritten{\psi}_2,x^*)\handwritten{\bigr\}}\handwritten{\unsure{\in|\epsilon}}t^{\handwritten{-r_i}}\cdot R.
\tag*{(A)}
\]
But $b_2^{-1}\handwritten{\psi}_2$ is an $R$-valued point of $G$, whose specialization, when $t\handwritten{\longrightarrow}0$, is the identity $e\handwritten{\in}G$; therefore $\handwritten{\sigma}^*(b_2^{-1}\handwritten{\psi}_2,x^*)$ is an $R$-valued point of $\handwritten{\mathbb{A}}^n$ whose specialization, when $t\handwritten{\longrightarrow}0$, is $x^*$; therefore
\[
X_i\handwritten{\bigl\{}\handwritten{\sigma}^*(b_2^{-1}\handwritten{\psi}_2,x^*)\handwritten{\bigr\}}=X_i(x^*)+t(Z_i)
\tag*{(B)}
\]
for some $Z_i\handwritten{\in}R$.

Combining (A) and (B), it follows that $r_i\handwritten{\geq}0$, whenever $X_i(x^*)\handwritten{\neq}0$. Therefore, $\handwritten{\mu}(x,b_2^{-1}\cdot\handwritten{\lambda}\cdot b_2)\handwritten{\leq}0$ and (i) is proven, by Proposition 2.3.
